\section{Análisis con los conceptos de la Unidad 1}

\subsection{Transparencias de distribución}

Netflix ofrece casi todas las transparencias. En el acceso, da igual si entro desde mi celular o la tele: todos consumen la misma API. Cumple con la ubicación porque, gracias a Zuul y Eureka, nadie necesita saber en qué servidor de AWS está el recurso~\cite{SpringCloudNetflix}. Ofrece transparencia de migración, ya que Netflix trasladó su sistema a AWS entre 2008 y 2016 sin que los usuarios hicieran nada~\cite{Netflix_CloudMigration}, y de replicación, porque cada servicio tiene varias copias y el mismo video está en muchos servidores de Open Connect, pero yo lo uso como si fuera uno solo~\cite{Netflix_OpenConnect}. También cumple la concurrencia, porque millones la usamos a la vez sin bloquearnos.

Lo que no ofrece del todo es la transparencia de reubicación y de fallos en el video: me ha pasado que la imagen se queda cargando un par de segundos, lo que ocurre cuando el servidor de Open Connect falla y el tráfico se redirige a otro~\cite{Netflix_OpenConnect}. El sistema se recupera, pero el fallo no fue invisible para mí.

\subsection{CAP y PACELC}

El teorema CAP dice que, ante una partición de red, un sistema distribuido debe elegir entre consistencia (datos iguales en todos los nodos) y disponibilidad (siempre responder, aunque el dato esté viejo)~\cite{vanSteen_2023}. Netflix prioriza la disponibilidad. Lo comprobé en mi experimento: tras la caída de una copia, Eureka siguió respondiendo unos 30 segundos con información desactualizada.

PACELC agrega qué pasa cuando no hay partición: se elige entre latencia y consistencia. Netflix prioriza la latencia. Por ejemplo, si termino un episodio en la laptop y abro la app en el celular enseguida, en \enquote{Seguir viendo} aún me sale por la mitad. Por eso Netflix se puede clasificar como PA/EL.

\subsection{Modelos de comunicación}

Netflix usa un modelo cliente-servidor: la app es el cliente y le pide todo a los servidores de Netflix. A diferencia de Spotify, que hasta 2014 usó un modelo híbrido con peer-to-peer~\cite{TechCrunch_Spotify2014}, Netflix no descarga contenido de otros usuarios. La comunicación hacia afuera es sincrónica: la app hace peticiones (REST o gRPC) y espera la respuesta~\cite{Rangarajan_PlayAPI}. Internamente también hay comunicación asincrónica, para tareas como guardar mi historial.

\subsection{Tolerancia a fallos}

Netflix está diseñado asumiendo que los servidores se van a caer. Esto viene de un incidente real: en 2008, una corrupción de su base de datos detuvo el envío de DVD durante tres días~\cite{Netflix_CloudMigration}.

Para detectar los fallos se basan en el monitoreo~\cite{Giamattei_2024}: los servicios mandan \enquote{latidos} a Eureka y, si dejan de llegar, el servicio se da por caído~\cite{SpringCloudNetflix}. En mi experimento, Eureka tardó unos 30 segundos en detectar la caída, mientras las otras dos copias seguían respondiendo. Para recuperarse, Hystrix \enquote{abre el circuito} para que el fallo no se propague~\cite{SpringCloudNetflix}. Además, con Chaos Monkey apagan servidores al azar para comprobar que el sistema se recupera solo~\cite{Owotogbe_2024}.
